% ============================================================================
\chapter{The Problem and the Promise}
\label{ch:problem}
% ============================================================================

\section{What This Book Is About}

This book describes a single, concrete neural network architecture and
explains every piece of it from the ground up. By the end, you will
understand not just \emph{what} the architecture does, but \emph{why}
each design choice was made and \emph{how} to implement it.

The architecture is a hybrid of three components:
\begin{enumerate}
  \item A \textbf{Tsetlin Machine}~\cite{granmo2018tsetlin} that learns interpretable Boolean
    rules from raw pixel data.
  \item A \textbf{Logic Gate Network}~\cite{petersen2022difflogic} that learns which logic gates
    best combine those rules into a classification decision.
  \item A \textbf{Lookup Table compiler}~\cite{liu2024kan} that converts any remaining
    learned functions into precomputed memory tables.
\end{enumerate}

The task is handwritten digit recognition on the MNIST dataset. The
implementation language is Rust. The ambition is simple: demonstrate
that a neural network built entirely from Boolean operations and memory
lookups can recognise digits with competitive accuracy---and do so
faster and with less energy than a conventional floating-point model~\cite{horowitz2014energy}.


\section{The MNIST Digit Recognition Task}
\label{sec:mnist-task}

The MNIST dataset~\cite{lecun1998mnist} is the ``Hello World'' of
machine learning. It consists of 70,000 images of handwritten digits
(0 through 9), each a $28 \times 28$ pixel grayscale image.

\begin{figure}[htbp]
\centering
\begin{tikzpicture}[scale=0.15]
  % Grid
  \fill[bitblue!5] (0,0) rectangle (28,28);
  \draw[bitgray!20, very thin, step=1] (0,0) grid (28,28);

  % Stylised "7"
  \foreach \x/\y in {
    8/24, 9/24, 10/24, 11/24, 12/24, 13/24, 14/24, 15/24, 16/24, 17/24, 18/24, 19/24,
    18/23, 19/23,
    17/22, 18/22,
    16/21, 17/21,
    15/20, 16/20,
    14/19, 15/19,
    14/18, 15/18,
    13/17, 14/17,
    13/16, 14/16,
    12/15, 13/15,
    12/14, 13/14,
    11/13, 12/13,
    11/12, 12/12,
    11/11, 12/11
  }{
    \fill[bitblue!60] (\x,\y) rectangle ++(1,1);
  }

  % Annotations
  \draw[<->, thick, bitgray] (-1, 0) -- (-1, 28)
    node[midway, left, font=\small\sffamily] {28\,px};
  \draw[<->, thick, bitgray] (0, -1) -- (28, -1)
    node[midway, below, font=\small\sffamily] {28\,px};
\end{tikzpicture}
\hspace{2cm}
\begin{tikzpicture}[scale=1]
  % Flattened vector
  \node[font=\sffamily\bfseries] at (0, 3.5) {Flattened};
  \draw[thick] (0, 0) -- (0, 3);
  \foreach \y/\v in {2.8/0, 2.5/0, 2.2/0.72, 1.9/0.95, 1.6/0.88,
    1.3/\ldots, 1.0/0, 0.7/0.31, 0.4/0} {
    \node[font=\tiny\ttfamily, right] at (0.1, \y) {\v};
  }
  \node[font=\scriptsize\sffamily, bitgray] at (0.5, -0.3)
    {784 values};

  % Arrow
  \draw[-{Stealth}, thick] (1.5, 1.5) -- (3, 1.5);
  \node[font=\scriptsize\sffamily] at (2.25, 1.9) {classify};

  % Output
  \node[draw, rounded corners, fill=bitgreen!20, minimum width=1.5cm,
    minimum height=0.8cm, font=\sffamily\bfseries] at (4, 1.5) {7};
\end{tikzpicture}
\caption{The MNIST task. A $28 \times 28$ grayscale image is flattened into
  a vector of 784 values (each in $[0, 255]$ or normalised to $[0, 1]$),
  and the network must output the correct digit label.}
\label{fig:mnist-task}
\end{figure}

\begin{itemize}
  \item \textbf{Training set:} 60,000 labelled images used to teach the
    model.
  \item \textbf{Test set:} 10,000 labelled images, never seen during
    training, used to measure accuracy.
  \item \textbf{Baseline:} A standard two-layer floating-point neural
    network achieves approximately 98\% accuracy on the test set.
\end{itemize}

MNIST is deliberately simple. We choose it not to set records but to
have a well-understood benchmark where we can focus entirely on the
architecture rather than on data engineering or scale.


\section{Why Bitwise? The Energy and Hardware Argument}
\label{sec:why-bitwise}

Every neural network in production today is built on a single
operation: \concept{multiply-accumulate} (MAC). Multiply a number by
a weight, add to a running total. Repeat billions of times.

This operation is expensive. \Cref{tab:energy-cost} shows the energy
cost of different operations in a modern chip.

\begin{table}[htbp]
\centering
\caption{Energy cost per operation in 45\,nm technology~\cite{horowitz2014energy}.}
\label{tab:energy-cost}
\begin{tabular}{@{}lrr@{}}
  \toprule
  \textbf{Operation} & \textbf{Energy (pJ)} & \textbf{vs.\ Logic Gate} \\
  \midrule
  Bitwise AND/OR/XOR  & 0.1  & $1\times$ \\
  Integer add (32-bit) & 0.9  & $9\times$ \\
  FP multiply (32-bit) & 3.7  & $37\times$ \\
  SRAM read            & 5.0  & $50\times$ \\
  DRAM read            & 640  & $6{,}400\times$ \\
  \bottomrule
\end{tabular}
\end{table}

\begin{keyinsight}
A floating-point multiplication costs 37 times more energy than a
bitwise logic operation~\cite{horowitz2014energy}. But moving data from main memory costs
6,400 times more. Any architecture that keeps its model small enough
to fit in cache---and uses logic operations instead of
multiplication---wins on both fronts.
\end{keyinsight}

There is also a market structure argument. Training and running large
floating-point models requires GPUs, a market dominated by a single
manufacturer. Bitwise operations run on every CPU ever made. An
architecture that needs only AND, OR, XOR, integer addition, and table
lookups can run on a microcontroller, an FPGA, or a 20-year-old laptop.


\section{The Hybrid Thesis}
\label{sec:thesis}

Our thesis is:

\begin{quote}
\itshape
A three-stage pipeline---Tsetlin Machine feature extraction, Logic
Gate Network composition, and Lookup Table deployment---can classify
MNIST digits with competitive accuracy ($>$97\%) while using
\emph{zero} floating-point operations at inference time.
\end{quote}

\begin{figure}[htbp]
\centering
\begin{tikzpicture}[
  stage/.style={draw, rounded corners=5pt, minimum width=3cm,
    minimum height=1.6cm, align=center, font=\small\sffamily,
    line width=0.8pt},
  arr/.style={-{Stealth[length=5pt]}, thick},
]
  \node[stage, fill=lightbg] (in) at (0, 0)
    {\textbf{Input}\\$28 \times 28$ pixels};
  \node[stage, fill=bitgreen!10, draw=bitgreen] (tm) at (4.5, 0)
    {\textbf{Tsetlin Machine}\\Boolean clauses};
  \node[stage, fill=bitorange!10, draw=bitorange] (lgn) at (9, 0)
    {\textbf{Logic Gates}\\Learned circuit};
  \node[stage, fill=bitpurple!10, draw=bitpurple] (lut) at (13.5, 0)
    {\textbf{LUT Scoring}\\Table lookups};

  \draw[arr] (in) -- (tm) node[midway, above, font=\scriptsize\sffamily]
    {booleanise};
  \draw[arr] (tm) -- (lgn) node[midway, above, font=\scriptsize\sffamily]
    {clause votes};
  \draw[arr] (lgn) -- (lut) node[midway, above, font=\scriptsize\sffamily]
    {gate outputs};

  % Training phases
  \node[font=\scriptsize\sffamily, bitgreen] at (4.5, -1.4)
    {Phase 1: automata};
  \node[font=\scriptsize\sffamily, bitorange] at (9, -1.4)
    {Phase 2: backprop};
  \node[font=\scriptsize\sffamily, bitpurple] at (13.5, -1.4)
    {Phase 3: compile};
\end{tikzpicture}
\caption{The three-stage hybrid architecture. Each stage is trained
  separately, and each uses a different learning mechanism. At inference
  time, the entire pipeline runs on Boolean operations, integer addition,
  and table lookups.}
\label{fig:hybrid-pipeline}
\end{figure}

The training happens in three phases:
\begin{description}[leftmargin=1em]
  \item[Phase 1 (Tsetlin):]
    Train Boolean clauses on booleanised pixel data using Tsetlin
    automata~\cite{tsetlin1961}. No floating-point arithmetic---only integer increment/decrement
    and random number generation.
  \item[Phase 2 (Logic Gates):]
    Freeze the clauses. Train a Logic Gate Network~\cite{petersen2024difflogic} on the clause outputs
    using gradient descent with a softmax relaxation over gate types.
    This phase uses floating-point arithmetic during training only.
  \item[Phase 3 (LUT Compilation):]
    Compile any remaining continuous scoring functions into integer
    lookup tables. This is a one-time offline step.
\end{description}

After Phase~3, the inference model contains no floating-point operations.


\section{What the Reader Needs}
\label{sec:prerequisites}

This book assumes:
\begin{itemize}
  \item \textbf{Linear algebra:} Vectors, matrices, dot products. You
    should be comfortable with $\bvec{y} = \bmat{W}\bvec{x} + \bvec{b}$.
  \item \textbf{Basic probability:} Random variables, expected value,
    conditional probability. You do not need measure theory.
  \item \textbf{Logic:} AND, OR, NOT, truth tables. Chapter~2 reviews
    this in detail.
  \item \textbf{Programming intuition:} You should be comfortable
    reading pseudocode. Rust experience is helpful for Chapter~14 but
    not required.
\end{itemize}

You do \emph{not} need prior experience with machine learning,
neural networks, or AI. Every concept is introduced from scratch.


\section{Notation}
\label{sec:notation}

We use the following conventions throughout:

\begin{table}[htbp]
\centering
\caption{Notation conventions used in this book.}
\label{tab:notation}
\begin{tabular}{@{}ll@{}}
  \toprule
  \textbf{Symbol} & \textbf{Meaning} \\
  \midrule
  $x, y, z$ & Scalars (lowercase italic) \\
  $\bvec{x}, \bvec{w}$ & Vectors (bold lowercase) \\
  $\bmat{W}$ & Matrices (bold uppercase) \\
  $\{0, 1\}$ & Binary / Boolean values \\
  $\{-1, +1\}$ & Bipolar binary values \\
  $\bar{x}$ & Boolean complement (NOT $x$) \\
  $\wedge, \vee$ & Logical AND, OR \\
  $\oplus$ & Logical XOR \\
  $|S|$ & Cardinality of set $S$ \\
  $[n]$ & The set $\{1, 2, \ldots, n\}$ \\
  $\popcount(\cdot)$ & Number of 1-bits in a binary word \\
  \bottomrule
\end{tabular}
\end{table}
